\begin{comment}
\begin{table*}
  \centering
  \begin{tabular}{@{}lccccccccccccc@{}}
    \toprule
    \multicolumn{1}{c}{} & \multicolumn{8}{c}{Scenes seen during online training} & \multicolumn{8}{c}{Scenes not seen during online training} \\
    \cmidrule(r){2-13}
    Method & Dunnottar Castle & Colosseum & Bannerman Castle & Pantheon & Christ the Redeemer & Statue of Liberty & Leaning Tower & Fushimi Castle \\
    \midrule
    Random Walk & todo & test\\
    SCONE-Entropy~\cite{guedon-nips22-scone-surface-coverage} & todo \\
    SCONE~\cite{guedon-nips22-scone-surface-coverage} & todo\\
    MACARONS (Ours) & todo\\
    \bottomrule
  \end{tabular}
  \caption{Results.   Ours is better.}
  \label{tab:coverage_auc}
\end{table*}
\end{comment}
%

% \begin{comment}
\begin{table*}
  \centering
  %\begin{subtable}{0.7\linewidth}
  {%\scriptsize
  \scalebox{0.8}{
  \begin{tabular}{@{}lcccc@{}}
    \toprule
    % \multicolumn{1}{c}{} & \multicolumn{4}{c}{Method} \\
    % \cmidrule(r){2-5}
    3D scene & Random Walk & SCONE-Entropy~\cite{guedon-nips22-scone-surface-coverage} & SCONE~\cite{guedon-nips22-scone-surface-coverage} & MACARONS (Ours) \\
    \midrule
    Dunnottar Castle 
    & 0.330 $\pm$ 0.106 
    & 0.381 $\pm$ 0.041 
    & 0.650 $\pm$ 0.093 
    & \textbf{0.784} $\pm$ 0.025 \\
    Colosseum 
    & 0.312 $\pm$ 0.098
    & 0.477 $\pm$ 0.023
    & 0.532 $\pm$ 0.032
    & \textbf{0.629} $\pm$ 0.011 \\
    Bannerman Castle 
    & 0.316 $\pm$ 0.106
    & 0.534 $\pm$ 0.012
    & 0.552 $\pm$ 0.020
    & \textbf{0.716} $\pm$ 0.029 \\
    Pantheon 
    & 0.206 $\pm$ 0.064
    & 0.338 $\pm$ 0.030
    & 0.401 $\pm$ 0.030
    & \textbf{0.499} $\pm$ 0.019 \\
    Christ the Redeemer 
    & 0.518 $\pm$ 0.058
    & 0.836 $\pm$ 0.028
    & 0.833 $\pm$ 0.037
    & \textbf{0.865} $\pm$ 0.037\\
    Statue of Liberty 
    & 0.296 $\pm$ 0.188
    & 0.673 $\pm$ 0.025
    & 0.695 $\pm$ 0.020 
    & \textbf{0.708 }$\pm$ 0.014\\
    Pisa Cathedral 
    & 0.327 $\pm$ 0.121
    & 0.431 $\pm$ 0.024
    & 0.555 $\pm$ 0.033
    & \textbf{0.661} $\pm$ 0.008 \\
    Fushimi Castle 
    & 0.513 $\pm$ 0.091
    & 0.791 $\pm$ 0.020
    & 0.806 $\pm$ 0.029
    & \textbf{0.842} $\pm$ 0.019 \\
    \midrule
    Alhambra Palace 
    & 0.361 $\pm$ 0.045
    & 0.425 $\pm$ 0.046
    & 0.528 $\pm$ 0.030
    & \textbf{0.634} $\pm$ 0.019\\
    Neuschwanstein Castle 
    & 0.396 $\pm$ 0.090
    & 0.491 $\pm$ 0.039
    & 0.662 $\pm$ 0.041
    & \textbf{0.760} $\pm$ 0.020 \\
    Eiffel Tower 
    & 0.417 $\pm$ 0.074
    & 0.702 $\pm$ 0.016
    & 0.753 $\pm$ 0.013
    & \textbf{0.789} $\pm$ 0.010 \\
    Manhattan Bridge 
    & 0.382 $\pm$ 0.092
    & 0.422 $\pm$ 0.055
    & 0.745 $\pm$ 0.069
    & \textbf{0.825} $\pm$ 0.034 \\
    \midrule
    Average on all scenes 
    & 0.364 
    & 0.542 
    & 0.643 
    & \textbf{0.726} \\
    \bottomrule
  \end{tabular}
  }}
  %\end{subtable}
  %\hfill
  %\begin{subtable}{0.23\linewidth}
  %  \includegraphics[width=\linewidth]{images/coverage_curves/all_path_planning_curve_macarons.png}
  %\end{subtable}
  \caption{{\bf AUCs of surface coverage on large 3D scenes.} All methods use perfect depth maps as input except for MACARONS, which takes RGB images as input. We follow \cite{guedon-nips22-scone-surface-coverage} and compute the area under the curve representing the evolution of the total surface coverage during exploration. The 8 scenes above the bar were seen by MACARONS during self-supervised training (but with different, random starting camera poses and trajectories), and the 4 scenes below the bar were not. Other methods are trained on ShapeNet~\cite{chang-15-shapenet} with 3D supervision. Even if it only uses RGB images, our model MACARONS is able to outperform the baselines in large environments since, contrary to other methods, its self-supervised online training strategy allows it to scale its learning process to any kind of environment.}
  \label{tab:coverage_auc}
\end{table*}
% \end{comment}

\begin{comment}
\begin{table*}
  \centering
  %\begin{subtable}{0.7\linewidth}
  {\scriptsize
  %\addtolength{\tabcolsep}{2pt}
  % \resizebox{0.8\width}{!}{
  \scalebox{0.84}{
  
  \begin{tabular}{@{}lccccccccccccc@{}}
    \toprule
    \multicolumn{1}{c}{} & \multicolumn{12}{c}{3D Scene} \\
    \cmidrule(r){2-14}
    \multicolumn{1}{c}{} & \multicolumn{8}{c}{Seen during online training} & \multicolumn{4}{c}{Not seen during online training} \\
    \cmidrule(r){2-9} \cmidrule(r){10-13}
     & Dunnottar & & Bannerman & & Christ the & Statue of & Pisa & Fushimi & Alhambra & Neuschwanstein & Eiffel & Manhattan & \\
     Method & Castle & Colosseum & Castle & Pantheon & Redeemer & Liberty & Cathedral & Castle & Palace & Castle & Tower & Bridge & Mean\\
     \midrule
     Random Walk & 0.330 & 0.312 & 0.316 & 0.206 & 0.518 & 0.296 & 0.327 & 0.513 & 0.361 & 0.396 & 0.417 & 0. 382 & 0.364 \\
     SCONE-Entropy~\cite{guedon-nips22-scone-surface-coverage} & 0.381 & 0.477 & 0.534 & 0.338 & 0.836 & 0.673 & 0.431 & 0.791 & 0.425 & 0.491 & 0.702 & 0.422 & 0.542 \\
     SCONE~\cite{guedon-nips22-scone-surface-coverage} & 0.650 & 0.532 & 0.552 & 0.401 & 0.833 & \textbf{0.695} & 0.555 & 0.806 & 0.528 & 0.662 & 0.753 & 0.745 & 0.643 \\
     MACARONS (Ours) & \textbf{0.764} & \textbf{0.618} & \textbf{0.722} & \textbf{0.516} & \textbf{0.878} & 0.671 & \textbf{0.633} & \textbf{0.828} & \textbf{0.603} & \textbf{0.759} & \textbf{0.786} & \textbf{0.853} & \textbf{0.719} \\
     \bottomrule
  \end{tabular}
  }}
  %\end{subtable}
  %\hfill
  %\begin{subtable}{0.23\linewidth}
  %  \includegraphics[width=\linewidth]{images/coverage_curves/all_path_planning_curve_macarons.png}
  %\end{subtable}
  \caption{{\bf AUCs of surface coverage in large 3D scenes.} All methods use perfect depth maps as input except for MACARONS, which takes RGB images as input. To compare the different methods, we follow \cite{guedon-nips22-scone-surface-coverage} and compute the area under the curve representing the evolution of the total surface coverage during exploration, after 100 NBV iterations. The surface coverage is computed using the ground truth meshes. AUCs are averaged on multiple trajectories in each scene: The same starting camera poses are used for each method, to allow for fair comparison. For this experience, all weights of MACARONS are frozen and we only perform inference computation. The model is considered to have been trained on a set of previous scenes. In this regard, the 8 scenes above the bar were seen during online training (but with different, random starting camera poses and trajectories), and the 4 scenes below the bar were not. Even if it only uses RGB images, our model MACARONS is able to outperform the baselines in large environments since, contrary to other methods, its self-supervised online training strategy allows it to scale its learning process to any kind of environment.}
  \label{tab:coverage_auc_test}
\end{table*}
\end{comment}

\begin{comment}
\begin{figure}[t]
  \centering
  %\fbox{\rule{0pt}{2in} \rule{0.9\linewidth}{0pt}}
   \includegraphics[width=0.5\linewidth]{images/coverage_curves/all_path_planning_curve_macarons.png}

   \caption{Example of caption.
   It is set in Roman so that mathematics (always set in Roman: $B \sin A = A \sin B$) may be included without an ugly clash.}
   \label{fig:onecol}
\end{figure}
\end{comment}